% ===== §9 Conclusions: at the boundaries =====
\section{Conclusions at the Boundaries of the Frameworks}
\label{sec:conclusions}

\noindent
A paper that has established the absence of a metalanguage, of an external standpoint, and of a framework that possesses the whole cannot terminate in a single synthesising conclusion, the commitment of which would constitute, in its terminal operation, the settling it has throughout excluded. It terminates as the prior paper terminated: not in a conclusion but in conclusions, plural, each stated from within a single framework, to that framework's boundary, in dialogue with the others and absorbed by none. What follows is not a summary but the statement, by each phase, of what it establishes and what it does not, and a concluding specification of the absence of a single conclusion, which is the paper's most adequate conclusion.

\subsection{From phenomenology}
\label{sec:concl-phenom}

Phenomenology establishes that there is a datum, prior to theory, in which two parties co-apprehend a nameless value generated between them, the co-apprehension being identically the generation. It specifies the datum's structure, the of-two simultaneity, the reflexive witness, the directedness upon an opening, the identity of apprehension and generation, and requires that every theory answer to this datum and not displace it. Phenomenology does not establish why the datum has this structure, or what must obtain, ontologically, for it to be possible; it describes, and transmits its description. Pressed beyond description into explanation, it reaches its boundary, at which it consigns the question to the ontology, having held the development accountable to a datum it describes but does not ground.

\subsection{From ontology}
\label{sec:concl-onto}

The ontology establishes that relation, grammar, and value are three aspects of one generative system, being, manifest dynamics, phenomenon, whose dynamics is reconfigured by its own products, and within which the subject is generated rather than presupposed. It establishes that this dissolves the circularity, that grammar is real and non-possessible, and that the system is unfinished in its constitution. The ontology does not establish the determinate form the reflexive structure assumes in any register; it supplies the abstract form and consigns its content to the disciplines into which the form breaks. And it does not, on pain of self-contradiction, complete itself: an ontology of a totality in generation does not deliver the completed totality, and its terminal statement is an exclusion, of substance, of the external standpoint, of completion, which consigns it beyond itself, to phases it does not occupy and a whole it does not possess.

\subsection{From the quantum-structural grammar}
\label{sec:concl-grammar}

The quantum-structural grammar establishes, in formal terms: that relational value is superposed potential and not determinate fact; that joint attention is participatory measurement, the manner of attending being the selection of the basis that fixes the space of admissible co-generation; that each crystallisation is a collapse conjoined with a reconfiguration of the space of the admissible; that the geometric phase tracked by the series is the formal descendant of the quantum phase; and that the grammar is non-enumerable in the epistemic sense that specifies the relation's non-possessibility. The grammar establishes nothing concerning physical substrate, having declined that thesis, and nothing concerning justice: it supplies the form of participatory generation and is silent as to whether a given generation is good, goodness not being a determination its formalism contains. It specifies the manner of value's coming-to-be; it does not specify the cost at which the coming-to-be proceeds. At that boundary it consigns the question to political economy, which it does not absorb and which does not absorb it.

\subsection{From the formal-modelling dilemma}
\label{sec:concl-modelling}

The examination of formal modelling establishes that the standard formal tools, the grammars of the Chomsky hierarchy and the fixed-law dynamical systems, are not inadequate in power to the reflexive value cycle, but are made adequate only by the introduction of a fixed higher-order structure, a meta-value-grammar, of which the reflexive appearances are determinate products; that this presupposition is the presupposition of a metalanguage and a big Other; that it incurs the regress Lacan arrests with the thesis that there is no Other of the Other; and that the one class of formal frameworks compatible with the denial of that Other is the background-independent, of which group field theory is the structural exemplar. The examination does not establish that a background-independent formalism for the value cycle can be constructed; it establishes only the location of the dilemma, that the formal modelling of the reflexive cycle either posits the fixed meta-value-grammar the series denies or seeks a formalism whose own ground is emergent. It specifies the form the question takes; it does not resolve it, and consigns it to future formal work and to the Value-Foam programme.

\subsection{From political economy}
\label{sec:concl-political}

Political economy establishes that the circulation of value reconfigures the subject that creates it, necessarily; that this reconfiguration is divided into a negative phase, alienation and the depleted subject and the producer colonised by its product, and a positive phase, self-realisation and the developed subject; that the division is marked by the sign of an independently definable reproduction surplus; that the subject is produced within the cycle it produces; and that the return to the creator is a spiral, a sameness bearing a phase, and not a return to the origin. Political economy does not establish the manner in which, in the intimate and non-measurable register, the value whose circulation it analyses is generated; for that it requires the participatory grammar it does not itself supply. And it does not certify, from its own resources, that a given return reaches its creator, that being a matter of practice and justice under actual asymmetry rather than of structural analysis. It specifies the moral direction of the reconfiguration; it does not guarantee the direction in a given case. At that point it reaches its boundary.

\subsection{From the account of symbolic justice}
\label{sec:concl-exploitation}

The account of symbolic justice establishes that the possession of a relation's generative structure is the symbolic correlate of the possession of the means of production; that its possessive modification, the legislation of the grammar from a claimed external position, is exploitation in the symbolic register, the appropriation of a jointly generated generativity to a single possessor and the reduction of the other to a legislated object; that the guidance of derivations without the legislation of the productions is the corresponding generative justice; and that the poem and the legislative reform of a language are the paradigms of the two. It establishes that symbolic exploitation and settling are one operation, the appropriation of the structure being the settling of the process, and that both extinguish the value they would possess. The account does not establish a procedure for the recognition of symbolic exploitation in a given case, the recognition being a matter of the practice and its two questions rather than of structural analysis; and it does not, on its own resources, distinguish guidance from possession in every instance, the distinction being one the parties must sustain rather than compute. It specifies the form of symbolic injustice and the criterion of its avoidance; it does not relieve the practitioner of the judgement. At that boundary it consigns the question to the practice and the ethics.

\subsection{From the practice and its ethics}
\label{sec:concl-praxis}

The practice establishes a disposition: to act not from an external position but from within; to cultivate the field rather than construct the grammar; to subtract rather than add, to preserve divergence, to sustain the opening without occupying it, to preserve the made language as poetic; to pose the two questions and, when they turn negative, to withdraw the compulsion. The ethics establishes that no finite subject guarantees its own goodness, the guarantor being the infinite reason that does not exist, and that this unguaranteeability constitutes the good as ethical, polyphonic, and unfinished. The practice and the ethics do not establish the determinate action in a given case; they yield a disposition and a pair of questions, not a procedure, and they decline, on principle, the rule that would relieve the practitioner of the unguaranteeable risk. Their terminal statement is not an instruction but a renunciation: of certainty, of the guarantee, of the position external to the cycle.

\subsection{The absence of a single conclusion}
\label{sec:no-conclusion}

These conclusions do not sum. Each framework, developed to its boundary, reaches a limit and consigns to the next a question the next takes up in a vocabulary the first does not possess; and the result of the paper obtains not in any one of them but in the dialogue among them, in the manner in which each determines what another does not reach, such that cognition develops through their exchange rather than settling into the possession of any single frame.

\begin{claim}[The absence of a single conclusion as the adequate conclusion]
There is no master conclusion, not in consequence of a failure to attain one, but because the totality treated is unfinished and no framework possesses it externally (\xref{sec:dialmat}). The conclusions are plural, partial, and mutually irreducible of necessity and not by default: a single synthesising conclusion would be a settling, the cashing of the whole into a single framework's terms, and thereby alienation in the register of cognition. The dialogue among the frameworks, each stated to its boundary, each consigning what it does not establish, is itself a good cycle, a polyphony recursively reproducing meaning, and is the sole form a conclusion may assume without controverting the development of the paper. That the paper does not conclude is not its limitation but its fidelity: it terminates as the truth it treats terminates, which is to say it does not terminate but is transmitted, spiralling, to be traversed again.
\end{claim}

\noindent
There is, accordingly, no metalanguage in which the joint import of the frameworks is statable, and the absence is not a deficiency but the result itself: the demonstration, effected once more, that the whole is not possessed but generated, not concluded but continued. The paper has stated only what is statable from within the phases, and has reached its boundary at each limit, consigning the dialogue among the limits as the nearest approximation to a whole that a finite and unfinished truth admits. What is not concludable admits of being witnessed; and the paper accordingly proceeds from conclusion to a closing return, which concludes nothing and records the point of departure.
